\section{System Design}
\subsection{Client}
A \texttt{Client} is a simple entity used to interact with the system from the outside. They are able to send Read and Write request to
a specific replica and receive a response containing the result of the operation. A \texttt{Client} keeps a timeout timer for every
request issued and will ignore responses received after it ends. A \texttt{Client} does \textit{not} receive a response for the result of a
correct write operation in the event that an election takes place or the contacted replica crashes while the request is being processed.

\subsection{Replica}
\subsubsection{Coordinator}
To implement the coordinator logic within our distributed system, we integrated all coordinator functionalities directly into the \texttt{Replica} class using dedicated methods. Since a coordinator is inherently also a replica, this architectural choice allows any replica to react to specific messages and seamlessly transition into the coordinator role if it is the next node to be elected. 
To keep all nodes informed about the identity of the current leader, each replica stores the coordinator's index in a local variable, which is subsequently utilized to forward incoming write requests to the correct destination. Additionally, a boolean flag is provided to verify whether the replica itself is currently acting as the coordinator. When this flag is true, it triggers specific coordinator-only operations, such as broadcasting periodic heartbeat signals.

Beyond heartbeats, the coordinator is responsible for critical synchronization duties, including:
\begin{itemize}
    \item Sending \texttt{WriteOK} messages to the replicas upon reaching a consensus quorum.
    \item Broadcasting synchronization messages immediately after a successful election to align the state of the network.
    \item Restoring updates that were left pending at the time of
        the previous coordinator's crash
\end{itemize}

\subsubsection{Update Protocol}
To ensure updates are correctly replicated across the system we use a protocol that is similar to 2PC, but with some significant differences
based on the assumptions we consider as true for this project.
% In particular, knowing that we can count on a strict majority of replicas
% always being correct and that nodes never recover from crashing were the two most impactful assumptions for our decisions.

The protocol works as follows:
\begin{enumerate}
    \item A replica receives a write request from a client and forwards it to the coordinator (this is true even if the coordinator itself
        is the one receiving the request). A timeout is set, waiting for the coordinator to send the update broadcast.
    \item The coordinator broadcasts the update to all replicas, including itself.
    \item Replicas receive the update and respond with an \texttt{UpdateACK} to the coordinator. A timeout is set, waiting for the
        \texttt{WRITEOK} message.
    \item The coordinator counts the ACKs until it has received enough to reach the quorum. It then broadcasts the \texttt{WRITEOK} message.
    \item Replicas receive the \texttt{WRITEOK} message and apply the update
\end{enumerate}

If at any point a timeout expires before the corresponding message is received then the replica holding said timeout will consider the
coordinator as crashed and will begin the election protocol. The election protocol details how they system recovers from partial updates.

\subsubsection{Election protocol}
The election protocol is triggered whenever a replica receives a \texttt{CoordinatorCrashed} message. Upon reception, the replica removes the failed leader from its list of available nodes and broadcasts an \texttt{ElectionStarted} message to initiate the protocol among the remaining peers. Replicas then remove the current coordinator from their active set and dynamically shift their operational state from what we call \texttt{Normal Mode} to \texttt{Election Mode}. 
We decoupled these two operational behaviors to handle concurrent user requests safely during an election phase. Consequently, any write request received while the system is in \texttt{Election Mode} is buffered locally and its forwarding is deferred until the new coordinator has been successfully established.
To optimize network bandwidth and ensure that ideally only one election message circulates at a time, the election initiation is staggered based on replica identifiers.

For every election message sent or forwarded, the sender expects an acknowledgment (ACK) from the recipient. If the ACK is not received within a specific timeout, the sender assumes the target replica has crashed and forwards the election message to the subsequent node in the ring.
If an ACK failure occurs during the second phase of the election (i.e., when the election message traverses the ring for the second time), the replica removes the history entry previously inserted by that crashed node during the first phase. This ensures that the coordinator receives an up-to-date topology of active nodes. Furthermore, if the crashed node happened to be the potential new coordinator, it is removed immediately, allowing the remaining replicas to locate an alternative coordinator without restarting the entire election protocol from scratch.

If the lowest-ID replica crashes or fails to initiate the election within this window, the next available node in the virtual ring will take over after its respective timeout expires, and so forth. If a node is delayed but does not crash, multiple election processes might concurrently circulate in the network; however, this does not compromise system correctness, as all replicas will eventually converge on the same coordinator.
The core of the election protocol strictly follows the project specifications, with an added optimization during the post-election phase. Once a new coordinator is elected, it reverts to \texttt{Normal Mode} and updates its view of active replicas based on which nodes contributed the latest updates during the election phase. 
The coordinator then broadcasts personalized synchronization messages tailored to each individual replica's missing updates.

When a non-coordinator replica detects that a coordinator can be elected but it is not itself the chosen candidate, it initializes a timeout proportional to the maximum number of replica in the system. If the expected synchronization message does not arrive within this window, the replica assumes the newly elected coordinator has crashed, thereby initiating a new election cycle.

Upon successfully receiving the \texttt{Synchronization} message, replicas initiate a recovery phase to retrieve any updates that may have been lost during the coordinator failure. Specifically, an update might have reached a quorum of replicas, yet the previous coordinator failed to broadcast the corresponding \texttt{WriteOK} message due to a sudden crash. To resolve such inconsistencies, each replica receiving the synchronization message responds by sending the new coordinator a list of all pending updates, those received from the previous coordinator for which no \texttt{WriteOK} confirmation was ever delivered.
The new coordinator evaluates these lists to determine whether any pending update has successfully satisfied the quorum requirement. If the quorum is confirmed, the new coordinator redistributes the update to all operational replicas, enabling them to safely apply it to their local state. Once this state transfer is complete, each replica switches its operational state back to \texttt{Normal Mode}, increments the system epoch counter, updates the active coordinator identifier, applies all missing updates, and clears its local queue of pending updates. This mechanism strictly enforces the main system property, guaranteeing that either every correct replica applies the update or none of them do.

Finally, crashes are simulated in strict accordance with the project specifications. The system evaluates whether a crash condition has been met and tracks the number of remaining messages before a crash occurs every time a message type supported by the crash configuration is processed.
